\documentclass[11pt,a4paper]{article}

% ---------- Layout ----------
\usepackage[margin=1in]{geometry}
\usepackage{microtype}

% ---------- Math ----------
\usepackage{amsmath,amssymb,amsthm,mathtools}

% ---------- Tables, links, bibliography ----------
\usepackage{booktabs}
\usepackage{graphicx}
\usepackage{hyperref}
\hypersetup{colorlinks=true,linkcolor=blue!60!black,urlcolor=blue!60!black,citecolor=blue!60!black}
\usepackage[round,authoryear]{natbib}

% ---------- Title ----------
\title{Learned Latent Curves and the Hyperspherical-Harmonic Variant:\\
  {\large A Closed-Loop Framework for Corpus Embedding and Recursive Skill Auditing}}

\author{Shant Tchatalbachian}

\date{August 6, 2026}

\begin{document}

\maketitle

\begin{abstract}
We describe a closed-loop framework for auditing and prioritizing improvement work across a corpus of engineering artifacts, and we document a sphere-aware extension of its curve-fitting stage together with the \emph{smallest audit unit} that the framework admits. The framework has four composed techniques: (1) a \emph{learned latent curve}, which re-expands a single 1-D ordering coordinate $t$ into a fixed-width $D$-dimensional embedding via a Fourier basis with learned frequencies; (2) \emph{curve-guided recursive self-improvement}, which uses sparse cells of the fitted curve as a prioritization lens for gap-mapping and bounded fixpoint-style editing cycles; (3) a \emph{hyperspherical-harmonic curve}, which replaces the flat $[0,1]^2$ parameter manifold of (1) with the Riemann sphere $S^2$ and a learned M\"obius $\phi_\theta \in \mathrm{PSL}(2,\mathbb{C})$ reparameterization; and (4) the \emph{single-action atom}, which is the smallest unit of the loop --- one corpus item maps to one point on $S^2$, one missing-primitive flip is one action, and the geodesic-only criterion gives an invariant that composes linearly across files. We give the governing equations for each technique, the empirical validation on a corpus of software skills ($N = 49$, $N = 70$, and $N = 79$ across three dated snapshots and one full-corpus audit), and a $474$-dispatch atom experiment on the $79$-skill corpus showing zero negative $\Delta$ and confirming Lemma~1 and Theorem~1 directly.
\end{abstract}

\section{Introduction}

The paper is organized so each main-body section earns its place by improving fit, clarifying the geometry, or sharpening the claim. Sections~2--4 give the family background (flat curve model, hyperspherical variant, atom) at the level of governing equations. Sections~5--6 give the dataset, evaluation protocol, and headline results, with the 79-skill corpus audit reported as the primary empirical exhibit. Sections~7--8 are scope and conclusion. Two appendices carry the remaining material.

\section{Background: The Learned-Latent-Curve Family}
\label{sec:family}

\subsection{Flat curve model}

For output dimension $j = 1, \dots, D$ (canonically $D = 384$) and 1-D coordinate $t \in [0, 1]$, the model is
\begin{equation}
z_j(t) = a_{j,0} + \sum_{m=1}^{k} \big( a_{j,m} \sin(2\pi f_m t) + b_{j,m} \cos(2\pi f_m t) \big),
\label{eq:flat}
\end{equation}
with $k$ shared learned frequencies $f_1, \dots, f_k$ and per-output coefficients $a_{j,m}, b_{j,m}$. Stacked over $j$, this is a curve $\gamma : \mathbb{R} \to \mathbb{R}^D$. Writing the design vector
\begin{equation}
\phi(t) = \big[\, 1,\; \sin(2\pi f_1 t),\, \cos(2\pi f_1 t),\, \dots,\, \sin(2\pi f_k t),\, \cos(2\pi f_k t) \,\big] \in \mathbb{R}^{1+2k},
\label{eq:design}
\end{equation}
the model is $\gamma(t) = C\, \phi(t)$ with coefficient matrix $C \in \mathbb{R}^{D \times (1+2k)}$, giving a parameter count of
\begin{equation}
P_{\mathrm{flat}} = D \cdot (1 + 2k).
\label{eq:paramflat}
\end{equation}
At $D = 384$, $k = 8$: $P_{\mathrm{flat}} = 384 \cdot 17 = 6{,}528$ parameters.

\subsection{Closed-form coefficient solve}

With the frequencies $f$ held fixed, the linear coefficient solve is the Tikhonov ridge,
\begin{equation}
C^{\star} = (\Phi^{\top}\Phi + \lambda I)^{-1} \Phi^{\top} Z,
\label{eq:ridge}
\end{equation}
where $\Phi \in \mathbb{R}^{N \times (1+2k)}$ stacks the design vector and $Z \in \mathbb{R}^{N \times D}$ stacks the target vectors. This is the sanity floor for any gradient-descent fit at the same frequencies: a fit worse than this at the same $f$ means the optimizer failed, not the model.

\subsection{Obtaining the coordinate $t$}
The default pipeline sets $t$ from the top principal component of a $z$-scored $N \times D_{\mathrm{feat}}$ feature matrix, mapped to $[0,1]$; a rank-uniformized variant replaces raw PC1 scores by their fractional ranks before mapping. The PC1+PC2 explained-variance ratio is the fit-quality gate ($\geq 0.40$); flat curves below the gate are not used.

\section{The Hyperspherical-Harmonic Curve}
\label{sec:variant}

\subsection{Model}

For $\mathbf{x} \in S^2 \subset \mathbb{R}^3$ (the Riemann sphere),
\begin{equation}
\gamma(\mathbf{x}) = \mathbf{b} + \sum_{\ell=0}^{L} \sum_{m} \mathbf{a}_{\ell,m}\, Y^{S^2}_{\ell,m}\!\big(\phi_\theta(\mathbf{x})\big),
\label{eq:hyperspherical}
\end{equation}
where $\mathbf{b} \in \mathbb{R}^D$ is the per-dim bias, $\mathbf{a}_{\ell,m} \in \mathbb{R}^D$ are the per-dim per-basis coefficients, and $\phi_\theta$ is the M\"obius reparameterization.

The parameter count is $D \cdot n_{\mathrm{basis}} + D + n_{\mathrm{M\"obius}}$, where $n_{\mathrm{basis}} = \sum_{\ell=0}^{L}(2\ell+1)$ and $n_{\mathrm{M\"obius}}$ is the real dimension of $\mathrm{PSL}(2,\mathbb{C})$ (Section~\ref{sec:mobius}). At $L=3$ on $S^2$, $n_{\mathrm{basis}}=16$; at $D=384$: $P_{\mathrm{sphere}} = 384 \cdot 16 + 384 + 6 = 6{,}534$.

\subsection{M\"obius reparameterization}
\label{sec:mobius}

$\phi_\theta(z) = (az+b)/(cz+d)$ with $a,b,c,d \in \mathbb{C}$ and $ad-bc=+1$. The four complex coefficients $a,b,c,d$ carry $8$ real components. The constraint $ad-bc=+1$ is a \emph{complex} constraint (real part equals 1, imaginary part equals 0), so it removes $2$ real degrees of freedom, leaving $6$ real degrees of freedom. The $\mathrm{PSL}(2,\mathbb{C})$ identification (matrices $M$ and $-M$ map to the same M\"obius transformation) is a discrete identification and does not further reduce the real dimension. So $\phi_\theta$ has $6$ real degrees of freedom, parameterised by the real and imaginary parts of $a,b,c,d$ subject to $ad-bc=1$.

We refine $\theta$ via L-BFGS-B with the closed-form ridge (Eq.~\ref{eq:ridge}) re-solved at each step. The basis is fixed; only the domain is reparameterized. Because every $\varphi \in \mathrm{PSL}(2,\mathbb{C})$ preserves the cross-ratio identically by definition, the cross-ratio check verifies implementation consistency: it fails if $\varphi_\theta$ is miscoded, not if the learned map is a poor fit.

\section{The Single-Action Atom and Linear Composition}
\label{sec:atom}

The smallest audit unit of the curve-guided framework is the \emph{single-action atom}. We define it precisely so the only-positive-$\Delta$ invariant it carries propagates linearly across the corpus.

\subsection{Atom: from file to $S^2$ point}

For a single corpus item $f$ (a file with $N_{\mathrm{sec}} \geq 2$ sections), the atom computes:
\begin{enumerate}
\item A \textbf{per-section coverage vector} $c_i \in \{0,1\}^9$, scored by pattern-matching against a fixed $9$-D binary primitive basis $\mathcal{P} = (p_0, \ldots, p_8)$.
\item A \textbf{weighted aggregate} $c = \sum_i w_i c_i$, with weights $w_i = \mathrm{len}(\mathrm{section}_i) / \mathrm{len}(f)$, thresholded at $0.5$.
\item A \textbf{section coverage matrix} $M \in \{0,1\}^{N_{\mathrm{sec}} \times 9}$, PCA top-$2$ via SVD giving $(\bar u, \bar v)$ for the file aggregate.
\item A \textbf{stereographic lift} $\sigma: \mathbb{R}^2 \to S^2$ with
\begin{equation}
\sigma(u,v) = \frac{1}{1+u^2+v^2}\big(2u,\; 2v,\; u^2+v^2-1\big) \in S^2,
\label{eq:stereographic}
\end{equation}
giving one point $p = \sigma(\bar u, \bar v) \in S^2$ per file.
\item An \textbf{ideal pole} $p^* = \sigma(\bar u^*, \bar v^*)$, where $(\bar u^*, \bar v^*)$ is the lift of the all-ones coverage vector $(1,\dots,1)$.
\item A \textbf{geodesic gap} $d(f) = \|p - p^*\|_2$ (chordal proxy on $S^2$).
\end{enumerate}
For each missing primitive $i \in \{j : c_j = 0\}$, the atom simulates the flip $c_i := 1$, recomputes the $S^2$ point $p'$, and selects $i^* = \arg\min_i d_{\mathrm{post}}(f)$. This is the \emph{geodesic-only criterion}: the chosen action strictly minimizes post-flip geodesic distance to the ideal pole.

\subsection{Lemma 1 (atom invariant)}

\begin{lemma}
For any file $f$ and any action $\alpha$ selected by the geodesic-only criterion, $\Delta_f = d_{\mathrm{pre}} - d_{\mathrm{post}} \geq 0$.
\end{lemma}

\begin{proof}
The criterion selects $\alpha^* = \arg\min_i d_{\mathrm{post}}$ over the $k$ candidates where $k = |\{j : c_j = 0\}|$. Each candidate flips exactly one missing primitive $i$ from $0$ to $1$ and recomputes the file's $S^2$ point. If all $k$ candidates had $d_{\mathrm{post}} \geq d_{\mathrm{pre}}$, the $\arg\min$ would tie at $d_{\mathrm{pre}}$ and $\Delta_f = 0$, never negative. A strictly negative $\Delta_f$ would require $d_{\mathrm{post}} > d_{\mathrm{pre}}$ for the $\arg\min$, contradicting the minimization. The action space is append-only (single-primitive flips); no candidate removes coverage.
\end{proof}

\subsection{Theorem 1 (linear composition)}

\begin{theorem}
For a corpus $C$ with $|C| = N$ files, every multi-file action $\alpha_{\mathrm{corpus}} = (\alpha_1, \ldots, \alpha_N)$ where each $\alpha_i$ is an atomic action on file $f_i$, has corpus-level
\begin{equation}
\Delta_{\mathrm{corpus}} = \sum_{i=1}^{N} \Delta_{f_i}.
\label{eq:composition}
\end{equation}
If every atomic $\Delta \geq 0$, then $\Delta_{\mathrm{corpus}} \geq 0$, and $\Delta_{\mathrm{corpus}} > 0$ if at least one atomic $\Delta > 0$.
\end{theorem}

\begin{proof}
Each $\alpha_i$ operates on its own file $f_i$ independently: the per-file coverage matrix and $S^2$ point are unchanged for all $f_j$ with $j \neq i$. The geodesic distance $d_{f_i}$ is a function of $f_i$'s coverage alone. Linear sum of non-negative scalars is non-negative.
\end{proof}

\subsection{M\"obius refinement strategy}

The corpus-level $\phi_\theta \in \mathrm{PSL}(2,\mathbb{C})$ is one M\"obius transformation applied uniformly to all files; the atom-bound composition rule (Theorem~1) requires it \emph{stable across per-cycle atom dispatches}, otherwise per-file $\Delta$s are measured in different charts and the linear sum does not apply to the reported cycles. Reported runs use the \emph{refine-once, $\phi_\theta$ frozen} mode ($\phi_\theta$ fit at corpus creation, frozen for all subsequent cycles); the joint-refine-per-cycle mode is reserved for $N_{\mathrm{items}} < 30$ or corpus growth $> 25\%$ since last refine.

\section{Dataset and Evaluation Protocol}
\label{sec:dataset}

\subsection{Dataset}

The corpus consists of software-skill specifications (engineering artifacts) in the yubiOS repository. Each item has a 9-dimensional binary feature vector recording which of nine primitive capabilities the skill implements. The corpus has grown across three dated snapshots: the alphabetical-first-half (49 items, $N_{\mathrm{train}}=35$, $N_{\mathrm{holdout}}=14$), the full directory at the time of the headline ablation (70 items, $N_{\mathrm{train}}=49$, $N_{\mathrm{holdout}}=21$), and the complete skill directory at this paper's revision (79 items, dated 2026-08-06, used for the corpus-audit RSI in Section~6.3 and the 474-dispatch atom experiment in Section~6.4).

The 9-D binary feature space has principal-component concentration $\mathrm{PC1}+\mathrm{PC2} = 0.652$ at the 49-item snapshot and $0.548$ at the 70-item snapshot; both clear the $\mathrm{PC1} \geq 0.40$ gate (Section~\ref{sec:family}).

\subsection{Baseline}

The capacity-matched baseline is a flat Fourier curve on $[0,1]^2$ with $k=2$ (the 2-D tensor-product extension of Eq.~\ref{eq:flat}), giving $5 \times 5 = 25$ basis functions and $9{,}984$ parameters. We choose $k=2$ because it is the lowest-order 2-D tensor product that exceeds the 16-function count of the spherical variant ($16$ basis, $6{,}534$ parameters); $k=1$ gives only $3 \times 3 = 9$ basis and would not be a capacity match for the 16-function sphere, while $k=3$ gives $7 \times 7 = 49$ basis and $\sim 24$k parameters, which overwhelms the corpus. Both models are fitted with the closed-form ridge (Eq.~\ref{eq:ridge}) and the same ridge regularisation $\lambda$.

\subsection{Evaluation metric}

The headline metric is \emph{matched-parameter ablation}: holdout $R^2$ at fewer or equal parameters, with the same split and the same ridge regularisation. We report the absolute holdout $R^2$ (not just the delta) so the result is interpretable against the corpus-mean baseline ($R^2 = 0$).

\subsection{Reproducibility}

All headline ablation numbers are single-run point estimates (no error bars) on a fixed holdout split, with a shared ridge regularisation $\lambda$ across both arms. The audit-phase numbers (Sections~6.3) carry 5-seed $\pm$ std error bars for phases E--H, where the 5-seed multi-cycle stress test was run; the headline ablation does not. All audit runs reported in this paper use the 79-skill corpus dated 2026-08-06 (single dated snapshot, referenced throughout).

\section{Results}

\subsection{Hyperspherical-harmonic variant: matched-parameter ablation}

The matched-parameter ablation on the two headline corpus snapshots is the paper's single contribution: on these corpora, the hyperspherical parameter manifold is a strictly better inductive bias than the flat $[0,1]^2$ baseline, with the absolute holdout $R^2$ positive on the smaller snapshot and the relative $\delta$ positive on both snapshots.

\textbf{49-item snapshot ($N_{\mathrm{train}}=35$, $N_{\mathrm{holdout}}=14$):}
\begin{itemize}
\item \textbf{Hyperspherical model} (S$^2$/$L{=}3$, 16 basis functions, 6,534 parameters): $R^2_{\mathrm{holdout}} = +0.618$.
\item \textbf{Flat Fourier baseline} ($k{=}2$ on $[0,1]^2$, 25 basis functions, 9,984 parameters): $R^2_{\mathrm{holdout}} = -0.359$.
\item \textbf{Matched-parameter $\delta$}: $+0.977$. The sphere wins by nearly 1.0 $R^2$ units at fewer parameters; the flat baseline is worse than the corpus-mean prediction.
\end{itemize}

\textbf{70-item snapshot ($N_{\mathrm{train}}=49$, $N_{\mathrm{holdout}}=21$, variant-included):}
\begin{itemize}
\item \textbf{Hyperspherical model}: $R^2_{\mathrm{holdout}} = +0.222$.
\item \textbf{Flat Fourier baseline}: $R^2_{\mathrm{holdout}} = -1.120$.
\item \textbf{Matched-parameter $\delta$}: $+1.342$. The sphere wins by more than 1.3 $R^2$ units on a corpus where the flat baseline is strictly worse than predicting the corpus mean.
\end{itemize}

\subsection{Corpus audit across cycles 5--9 (phases A--H)}

\begin{figure}[h]
\centering
\includegraphics[width=5.5in]{chart-A-H-1-progression.png}
\caption{\textbf{Figure 1.} Phase A $\rightarrow$ B $\rightarrow$ C $\rightarrow$ D $\rightarrow$ E $\rightarrow$ F $\rightarrow$ G $\rightarrow$ H holdout $R^2$ progression (cycles 5--9 RSI corpus audit, 70$\rightarrow$73-skill corpus). Error bars = 5-seed $\pm$ std for E, F, G, H. The sphere arm climbs $-0.5021 \to +0.2534$ while the flat arm stays flat at $-0.4588 \to -0.4231$; the win is not a single-split artifact.}
\label{fig:progression}
\end{figure}

\begin{table}[h]
\centering
\includegraphics[width=5.5in]{chart-A-H-table-1-headline.png}
\caption{\textbf{Table 1.} Headline numbers across cycles 5--9 (phases A--H; 70$\rightarrow$73-skill corpus).}
\label{tab:headline}
\end{table}

\subsection{Atom experiment: 474 dispatches on the 79-skill corpus}

The 79-skill corpus was audited via $474$ atom dispatches across $6$ cycles (single run on the 2026-08-06 snapshot, commit $\texttt{6ae3abeb65}$ on `yubi-OS/yubiOS main`). Cumulative $\Delta = +11.2963$ across all dispatches (sum of the per-cycle aggregates; the unrounded sum over all $474$ individual dispatches is $+11.2971$); zero negative $\Delta$; $116$ strictly positive $\Delta$ ($358$ saturated files contributed $\Delta = 0$, see Appendix~\ref{app:atom-coverage}). Sparse cells: $7 \to 3 \to 2 \to 3 \to 0 \to 0$. Fixpoint reached at cycle 6 (peak $\Delta$ below the $0.001$ epsilon). The corpus reached $79/79 = 100\%$ coverage across all nine primitives. The $20$-cycle $11$-file deep-research experiment is in Appendix~\ref{app:20cycle}.

\begin{figure}[h]
\centering
\includegraphics[width=6.5in]{chart-RSI-79-cycles-1-6.png}
\caption{\textbf{Figure 2.} 79-skill corpus RSI: cumulative $\Delta$ per cycle (left) and per-primitive coverage progression (right). Six cycles to fixpoint. Total: 474 atom dispatches, 0 negative $\Delta$, cumulative $\Delta = +11.2963$ (sum of per-cycle aggregates; the unrounded sum is $+11.2971$). Per-cycle cumulative trajectory: $+5.5787 \to +3.0091 \to +1.2833 \to +0.8829 \to +0.5423 \to +0.0000$ --- the diminishing-returns curve predicted by single-action-curve-rsi's fixpoint rule.}
\label{fig:rsi-79}
\end{figure}


\section{Discussion}

\subsection{What this result does and does not show}

The hyperspherical-harmonic variant wins on the matched-parameter ablation at both corpus snapshots by a margin (49-item $\delta = +0.977$, 70-item $\delta = +1.342$) that is hard to attribute to noise. We do \emph{not} claim the variant is a strict improvement over the flat curve in absolute terms: on the 70-item snapshot the absolute $R^2$ is $+0.222$ (positive but small), and the headline numbers are single-run point estimates without error bars. The atom experiment (Figure~2) shows the smallest audit unit composes without regression --- 474 dispatches, 0 negative $\Delta$ --- but does not by itself validate the variant.

\subsection{Limitations}

The empirical case is fragile because it is one corpus family, one target, and the headline numbers are single-run point estimates. The 5-seed error bars in Figure~1 (phases E--H) partially address the empirical-fragility concern for the audit; the headline ablation numbers do not. The multi-corpus audit (Appendix~\ref{app:multi-corpus}) uses $1391$ dispatches across $213$ files ($79$ skills, $21$ docs, $113$ refs) --- enough to support an empirical claim at the atom-invariant level (zero negative $\Delta$ across all $1391$), but not enough to support multi-seed error bars at the headline-ablation level. A synthetic-manifold benchmark (Appendix~\ref{app:20cycle}.3, not run) and a second-corpus re-run of the ablation itself would replace the present single-seed point estimates with error bars at both levels.

\section{Conclusion}

The hyperspherical-harmonic curve replaces the flat $[0,1]^2$ parameter manifold with the Riemann sphere $S^2$ and learns a M\"obius reparameterization $\phi_\theta \in \mathrm{PSL}(2,\mathbb{C})$ of the domain. On the yubiOS software-skill corpus, it wins on the matched-parameter ablation at both dated snapshots (49-item $\delta = +0.977$, 70-item $\delta = +1.342$), with fewer parameters and no error bars. The single-action atom is the smallest unit of the resulting audit pipeline; its only-positive-$\Delta$ invariant propagates linearly to multi-file composition (Theorem~1), as confirmed by $1391$ atom dispatches across the three corpora (\texttt{skills/}, \texttt{docs/}, \texttt{refs/}) in Appendix~\ref{app:multi-corpus} and a $20$-cycle experiment on the $11$-file deep-research corpus in Appendix~\ref{app:20cycle}, showing zero negative $\Delta$ across all of them. Future work: a synthetic-manifold benchmark (Appendix~\ref{app:20cycle}.3) would test whether the inductive-bias claim holds on manifolds other than $S^2$, and a second-corpus re-run of the ablation itself would replace the present single-seed point estimates with error bars.

\appendix

\section{Atom Coverage of 79 Skills (Empirical)}
\label{app:atom-coverage}

This appendix gives the per-cycle accounting behind the headline atom result in Section~6.3. The $79$-skill corpus (\texttt{skills/} on \texttt{yubi-OS/yubiOS main}, snapshot 2026-08-06) was audited by $474$ atom dispatches --- $79$ files $\times$ $6$ cycles, one dispatch per file per cycle. Of those $474$ dispatches, \textbf{$116$ produced a strictly positive $\Delta$}, \textbf{$358$ produced $\Delta = 0$} (the file was already saturated on all nine primitives, so the atom had no candidate flip), and \textbf{zero produced a negative $\Delta$}. The zero-$\Delta$ majority is the expected shape, not a defect: Lemma~1 guarantees non-negativity, and a saturated file has an empty candidate set, so the arg-min ties at $d_{\mathrm{pre}}$. Cumulative $\Delta$ = $\mathbf{+11.2963}$ (sum of the per-cycle aggregates as published in \texttt{rsi-3-corpus-summary-2026-08-06.json}; the unrounded sum over all $474$ individual dispatches is $+11.2971$).

\subsection{Per-cycle trajectory}

Per-cycle cumulative $\Delta$: $+5.5787 \to +3.0091 \to +1.2833 \to +0.8829 \to +0.5423 \to +0.0000$. Per-cycle count of strictly positive dispatches: $59 \to 30 \to 12 \to 9 \to 6 \to 0$. Sparse cells: $7 \to 3 \to 2 \to 3 \to 0 \to 0$. Both the $\Delta$ series and the positive-dispatch count fall monotonically; the sparse-cell count does \emph{not} ($3 \to 2 \to 3$ at cycles 3--4). This is expected: a single-primitive flip moves a file's $S^2$ point, so a cycle can vacate one equal-area cell and isolate a neighbour. Lemma~1 protects $\Delta$, not the sparse-cell count. Fixpoint was declared at cycle 6, where peak $\Delta = 0.0000$ falls below the $0.001$ epsilon.

\subsection{Per-primitive coverage progression}

Coverage counts out of 79, ordered (attestation, trust\_chain, least\_privilege, declarative\_policy, continuous\_adaptive, immutability, audit\_evidence, cryptographic\_identity, segmentation), recorded at the head of each cycle:
\begin{align*}
\mathbf{c_1} &= [71,\, 63,\, 63,\, 49,\, 35,\, 74,\, 72,\, 70,\, 73] \\
\mathbf{c_2} &= [79,\, 71,\, 69,\, 75,\, 54,\, 78,\, 77,\, 70,\, 73] \\
\mathbf{c_3} &= [79,\, 79,\, 71,\, 78,\, 70,\, 78,\, 77,\, 71,\, 75] \\
\mathbf{c_4} &= [79,\, 79,\, 79,\, 79,\, 72,\, 79,\, 78,\, 71,\, 79] \\
\mathbf{c_5} &= [79,\, 79,\, 79,\, 79,\, 79,\, 79,\, 79,\, 73,\, 79] \\
\mathbf{c_6} &= [79,\, 79,\, 79,\, 79,\, 79,\, 79,\, 79,\, 79,\, 79]
\end{align*}
\texttt{continuous\_adaptive} begins at $35/79$ --- the sparsest column in the corpus --- but the geodesic-only criterion targets it aggressively ($44$ of the $116$ winning flips) and it is saturated by cycle~5. \texttt{cryptographic\_identity} begins at a comfortable $70/79$ and is the last column to close, still open entering cycle~6. The criterion minimizes post-flip geodesic distance, not per-column deficit.

\subsection{Winner distribution and the quantization of $\Delta$}

Across the $116$ winning flips the selected primitive was \texttt{continuous\_adaptive} $44\times$, \texttt{declarative\_policy} $22\times$, \texttt{trust\_chain} $16\times$, \texttt{least\_privilege} $16\times$, \texttt{attestation} $8\times$, \texttt{cryptographic\_identity} $8\times$, \texttt{immutability} $1\times$, and \texttt{audit\_evidence} $1\times$. The realized $\Delta$ is a function of $k = |\{j : c_j = 0\}|$, the count of missing primitives, and of nothing else --- not of which primitive is flipped, and not of which file. The observed map is
\begin{align*}
k = 1 &\to 0.0904, \quad 2 \to 0.1003, \quad 3 \to 0.1098, \\
k = 4 &\to 0.1175, \quad 5 \to 0.1206, \quad 6 \to 0.1158, \\
k = 7 &\to 0.0989, \quad 8 \to 0.0678, \quad 9 \to 0.0242.
\end{align*}
The map is unimodal with its maximum at $k = 5$: a file missing five of nine primitives sits where the chordal metric on $S^2$ is steepest. The maximum single-dispatch $\Delta$ in the run was $0.1206$, first attained at cycle~1 by \texttt{negative-skill-space} ($k = 5$, winner \texttt{trust\_chain}) and \texttt{single-action-curve-rsi} ($k = 5$, winner \texttt{attestation}). The identical ladder is reproduced on the \texttt{docs/} and \texttt{refs/} corpora (Appendix~\ref{app:multi-corpus}).

\begin{figure}[h]
\centering
\includegraphics[width=5.5in]{chart-A-H-2-per-cycle-delta.png}
\caption{\textbf{Figure A.1.} Per-cycle $\Delta$ across the phases A--H corpus audit. The diminishing-returns envelope is the operational signature of the fixpoint rule: each cycle closes the highest-$\Delta$ candidates, leaving a strictly cheaper residual for the next.}
\label{fig:app-a1}
\end{figure}

\begin{figure}[h]
\centering
\includegraphics[width=5.5in]{chart-A-H-3-primitive-delta.png}
\caption{\textbf{Figure A.2.} Per-primitive $\Delta$ contribution. \texttt{continuous\_adaptive} dominates the total ($44$ of $116$ winning flips), consistent with its $35/79$ starting coverage.}
\label{fig:app-a2}
\end{figure}

\begin{figure}[h]
\centering
\includegraphics[width=5.5in]{chart-A-H-table-2-cycle-summary.png}
\caption{\textbf{Table A.1.} Per-cycle summary of the phases A--H audit (holdout $R^2$, both arms, 5-seed $\pm$ std where available).}
\label{tab:app-a1}
\end{figure}

\begin{figure}[h]
\centering
\includegraphics[width=5.5in]{chart-A-H-table-3-primitive-progression.png}
\caption{\textbf{Table A.2.} Per-primitive coverage progression across cycles, tabular form of the $c_1$--$c_6$ vectors in Section~A.2.}
\label{tab:app-a2}
\end{figure}

\section{Multi-Corpus RSI Audit --- \texttt{skills/}, \texttt{docs/}, \texttt{refs/}}
\label{app:multi-corpus}

The $79$-skill audit reported in Section~6.3 used cycles 10--15 of the \texttt{curve-guided-rsi} self pipeline; cycles 1--9 were the prior phases A--H audit on the $70$-skill corpus. That audit has since been extended from the engineering corpus to the whole repository. This appendix reports the multi-corpus run: the same single-action-atom pipeline, the same 9-D \texttt{internal-big-picture} primitive basis, the same chordal proxy on $S^2$ with $\phi_\theta$ frozen at identity, applied independently to the self-documentation corpus (\texttt{docs/}, $21$ files) and the references corpus (\texttt{refs/}, $113$ files) alongside the engineering corpus (\texttt{skills/}, $79$ files). All three reach corpus-level fixpoint.

\subsection{Headline totals}

\textbf{$1391$ atom dispatches across $213$ files, $598$ strictly positive $\Delta$, zero negative $\Delta$, cumulative $\Delta = +59.7671$.} By corpus:
\begin{itemize}
\item \texttt{skills/} --- $79$ files, $6$ cycles, $474$ dispatches, $116$ positive, $\Delta = +11.2963$, sparse cells $7 \to 0$.
\item \texttt{docs/} --- $21$ files, $6$ cycles, $126$ dispatches, $56$ positive, $\Delta = +5.5410$, sparse cells $7 \to 0$.
\item \texttt{refs/} --- $113$ files, $7$ cycles, $791$ dispatches, $426$ positive, $\Delta = +42.9298$, sparse cells $33 \to 0$.
\end{itemize}
All three corpora terminate on the same rule (peak $\Delta < 0.001$) and all three reach $100\%$ coverage on all nine primitives. This is the verification of Lemma~1 at scale: $1391$ independent geodesic arg-min selections, not one of which produced a regression. Source: \texttt{papers/data/rsi-3-corpus-summary-2026-08-06.json}.

\subsection{Differential curve baselines}

Raw cumulative $\Delta$ is not comparable across corpora of different size, so the differential baseline normalizes by corpus size: mean $\Delta$ per file per cycle. \texttt{skills/} opens at $0.0706$ $\Delta$/file and decays $0.0381, 0.0162, 0.0112, 0.0069, 0$. \texttt{docs/} opens at $0.0929$ and decays $0.0795, 0.0604, 0.0225, 0.0086, 0$. \texttt{refs/} opens at $0.0881$, \emph{rises} to $0.1061$ at cycle 2, then decays $0.0918, 0.0657, 0.0259, 0.0024, 0$. The \texttt{refs/} corpus is the one departure from the monotone diminishing-returns shape seen elsewhere in this paper: per-cycle $\Delta$ on \texttt{refs/} goes $+9.9559 \to +11.9840 \to +10.3710 \to +7.4256 \to +2.9221 \to +0.2712 \to 0$, peaking at cycle~2. Corollary~1 constrains the \emph{cumulative} sum to be monotone non-decreasing, which it is; it says nothing about the per-cycle increment. The mechanism is the $k$-quantization of Section~A.3: \texttt{refs/} begins with most files deep in the high-$k$ tail ($k = 7$--$9$, where $\Delta$ is small), cycle~1 walks them up into the $k \asymp 5$ band where $\Delta$ is maximal, and cycle~2 therefore harvests more total distance than cycle~1 did.

\subsection{What transfers across corpora}

Three properties hold identically on all three corpora and are therefore properties of the atom rather than of the engineering corpus. \textbf{First}, the $\Delta$ ladder: the realized $\Delta$ depends only on $k$, the count of missing primitives, with the same nine values and the same unimodal peak at $k = 5$, on \texttt{skills/}, \texttt{docs/} and \texttt{refs/} alike. \textbf{Second}, the peak-$\Delta$ termination sequence: all three corpora walk down the same ladder ($0.1206 \to \ldots \to 0.0904 \to 0$) as the highest-$\Delta$ band is exhausted. \textbf{Third}, the identity of the laggard primitive: \texttt{cryptographic\_identity} is the last column to saturate in every corpus. The geodesic-only criterion has the least incentive to act on a primitive missing from files already close to the ideal pole.

\subsection{Scope of the multi-corpus claim}

The multi-corpus audit strengthens the atom-invariant claim (Lemma~1 and Theorem~1 now hold over $1391$ dispatches on three structurally distinct corpora rather than $474$ on one) but it does \emph{not} strengthen the headline matched-parameter ablation of Section~6.1, which remains a single-seed point estimate on one corpus family. The second-corpus re-run called for in Section~7.2 would have to re-run the \emph{ablation}, not the audit, and has not been done.

\begin{figure}[h]
\centering
\includegraphics[width=4.8in]{chart-Appendix-B-3-corpus-2026-08-06.png}
\caption{\textbf{Figure B.1.} Multi-corpus hyper-sphere RSI audit. Top: cumulative $\Delta$ per cycle, raw, by corpus. Middle: the differential curve baseline --- $\Delta$ per file per cycle, normalized by corpus size. Bottom: sparse-cell count per cycle ($S^2$ equal-area partition, $r = 0.05$). \texttt{skills/} 6 cycles $\Delta = +11.2963$; \texttt{docs/} 6 cycles $\Delta = +5.5410$; \texttt{refs/} 7 cycles $\Delta = +42.9298$. Zero negative $\Delta$ across all $1391$ dispatches.}
\label{fig:app-b1}
\end{figure}

\section{The 20-Cycle Deep-Research Corpus and the Synthetic-Manifold Benchmark (Partly Future Work)}
\label{app:20cycle}

Sections~6 and Appendix~\ref{app:multi-corpus} report corpus-scale audits in which each file receives at most one dispatch per cycle. A complementary experiment runs the atom repeatedly against a small corpus to expose the per-file convergence behaviour that a wide, shallow audit averages away. This appendix reports that experiment and then sketches the one benchmark this paper still has not run.

\subsection{$20$ cycles on $11$ deep-research files}

The $11$-file deep-research corpus (the \texttt{sealed-uki-vm} and \texttt{falco} outputs in the yubiOS references tree) was audited over $20$ atom cycles --- an initial sweep plus seven post-edit re-fits --- with the corpus re-fitted after every applied edit, so each cycle sees a genuinely changed manifold. Zero cycles produced a negative $\Delta$. Cumulative $\Delta$ plateaued at $+1.6882$. Source: \texttt{papers/data/single-action-curve-rsi-cycles-2026-08-05.json}.

\subsection{Diminishing marginal value and the shifting peak}

The three peak runs trace $+0.3092$ (cycle~2, advisor-report, target \texttt{has\_source}) $\to +0.2705$ (the same file re-selected after its first edit, with a smaller $\Delta$ because the edit moved it closer to the ideal pole) $\to +0.1872$ (cycle~14, falco). Peak $\Delta$ fell $39.5\%$ across the three peak runs and mean $\Delta$ fell $28.2\%$, from $+0.1176$ to $+0.0844$. Per-file $\Delta$ reductions after editing: advisor-report $-55.7\%$, pkcs11-ecdsa-deepdive $-66.6\%$, pkcs11-ecdsa-VERIFIED $-47.5\%$, prior-art-V52 $-43.4\%$, comparative-report $-100\%$. Four of the eleven files reached a local minimum ($\Delta = 0$) and required no further action, up from two at cycle~12. The atom sweeps the corpus rather than over-fitting one file --- the same mechanism that produces the cycle-2 rise on \texttt{refs/} in Section~B.2.

\begin{figure}[h]
\centering
\includegraphics[width=6.5in]{chart-D-1-20-cycle-delta.png}
\caption{\textbf{Figure C.1.} $\Delta$ per cycle across the 20-cycle atom experiment on 11 deep-research files. Blue: initial corpus sweep (cycles 1--12). Red: post-edit re-fits (cycles 13--20). Stars mark the three peak runs. Cumulative $\Delta = +1.6882$; zero negative $\Delta$ across all 20 cycles.}
\label{fig:app-c1}
\end{figure}

\subsection{Synthetic-manifold benchmark (not run)}

Every empirical result in this paper is measured on corpora that were \emph{assumed} to suit an $S^2$ parameter manifold. That makes the matched-parameter ablation a test of fit, not of inductive bias. The benchmark that would separate the two is a negative control. \textbf{Corpus:} $N = 200$ synthetic points on the torus $T^2 = S^1 \times S^1$, with per-point 9-D binary feature vectors encoding toroidal modes. \textbf{Arms:} the hyperspherical-harmonic curve on $S^2$ ($L = 3$, $16$ basis functions) against the flat Fourier baseline on $[0,1]^2$ ($k = 2$, $25$ basis functions), both fitted by the closed-form ridge of Eq.~(4) at the same $\lambda$, on the same holdout split, over $5$ seeds. \textbf{Prediction:} the flat baseline wins on $T^2$ --- a sphere is the wrong prior for a genus-1 manifold --- while the sphere wins on an $S^2$-generated positive control. \textbf{Status: not implemented.} This benchmark and a second-corpus re-run of the ablation itself remain the two open items that bound the headline claim.

\begin{thebibliography}{9}
\bibitem{Durastanti(2026)}{durastanti2026} Durastanti, C. (2026). \emph{Spectral Bayesian Regression on the Sphere.} arXiv:2601.20528 [math.ST]. \url{https://arxiv.org/abs/2601.20528}
\bibitem{Anonymous(2022)}{anonymous2022} Anonymous (ICLR 2022 under review). \emph{Generalized Fourier Features for Coordinate-Based Learning on Manifolds.} OpenReview \texttt{g6UqpVislvH}.
\bibitem{Rahimi and Recht(2007)}{rahimi2007} Rahimi, A., \& Recht, B. (2007). \emph{Random Features for Large-Scale Kernel Machines.} NeurIPS 2007.
\bibitem{Tancik et~al.(2020)}{tancik2020} Tancik, M., et al. (2020). \emph{Fourier Features Let Networks Learn High Frequency Functions in Low Dimensional Domains.} NeurIPS 2020.
\bibitem{Sitzmann et~al.(2020)}{sitzmann2020} Sitzmann, V., Martel, J., Bergman, A., Lindell, D., \& Wetzstein, G. (2020). \emph{Implicit Neural Representations with Periodic Activation Functions.} NeurIPS 2020 (SIREN).
\bibitem{Mildenhall et~al.(2020)}{mildenhall2020} Mildenhall, B., et al. (2020). \emph{NeRF: Representing Scenes as Neural Fields for View Synthesis.} ECCV 2020.
\bibitem{Cohen et~al.(2018)}{cohen2018} Cohen, T. S., et al. (2018). \emph{Spherical CNNs.} ICLR 2018.
\bibitem{Nickel and Kiela(2017)}{nickel2017} Nickel, M. \& Kiela, D. (2017). \emph{Poincar\'e Embeddings for Learning Hierarchical Representations.} NeurIPS 2017.
\bibitem{Ahlfors(1979)}{ahlfors1979} Ahlfors, L. V. (1979). \emph{Complex Analysis,} 3rd~ed. McGraw-Hill.
\bibitem{do~Carmo(1976)}{docarmo1976} do Carmo, M. P. (1976). \emph{Differential Geometry of Curves and Surfaces.} Prentice-Hall.
\bibitem{Smith(2026)}{smith2026} Smith, E. (2026). \emph{Learning in Curved Weight Space: Exponential-Linear Weight Reparameterization for Improved Optimization.} arXiv:2607.09967 [cs.LG].
\end{thebibliography}

\end{document}
